
% #### 2. Nástroj 0: Nativní UI GitHubu
% 1. **Záložka Insights (Základní analytika):**
%    * *Pulse:* Poskytuje rychlý přehled (snapshot) o tom, co se dělo za poslední den/týden/měsíc (kolik se sloučilo PR, kolik se otevřelo/zavřelo Issues, nejaktivnější autoři).
%    * *Contributors & Commits:* Grafy ukazující historii aktivity. Kdo napsal kolik řádků, kdy byly největší špičky vývoje.
%    * *Traffic & Dependency graph:* Kdo repozitář klonuje a na jakých externích knihovnách projekt závisí.

% 2. **Správa Issues a Pull Requestů:**
%    * Možnost filtrovat pomocí vyhledávacího jazyka GitHubu (např. `is:pr is:open label:bug no:assignee`).
%    * *Milestones (Milníky):* Nativní způsob, jak sledovat postup prací na určité verzi (procento dokončení podle zavřených Issues).

% 3. **GitHub Projects (Nástěnky a tabulky):**
%    * GitHub nabízí integrovaný projektový management (Kanban boardy, tabulkové pohledy), kde lze práci vizualizovat a přidávat vlastní pole (custom fields).

% 4. **Záložky Actions a Security:**
%    * *Actions:* Nativní výpis běhu CI/CD pipelin. Zelené/červené ikonky ukazující, zda build prošel.
%    * *Security:* Integrovaný Dependabot a Code Scanning upozorňující na zranitelnosti v kódu a zastaralé závislosti.

\section{Nativní statistiky GitHubu}
\label{sec:github_stats}

GitHub nabízí základní nativní statistiky dostupné přímo u~každého projektu. Tyto statistiky jsou k~dispozici v~záložce \textit{Insights} a~zahrnují přehled otevřených a~zavřených \PR, aktivitu přispěvatelů a~další údaje.

\subsection{Obsah a~funkce}
\label{subsec:github_stats:content}

Hlavním místem pro získání obecných statistik a~grafů je záložka Insights. Nejdůležitější části zahrnují pohledy Pulse, Contributors a~Commits. U~všech statistik lze vybrat časové období, nejčastěji den, týden, měsíc nebo rok. U~některých statistik je z~důvodu velkého množství dat toto období omezeno, například na maximálně poslední měsíc.

\begin{description}
  \item[Pulse] Poskytuje rychlý přehled o~aktivních a~uzavřených \PR, otevřených a~uzavřených issues, nejaktivnějších autorech a~základní sumarizaci aktivit včetně přidaných a~odebraných řádků kódu. Obsahuje i~seznam uzavřených \PR za zvolené časové období. Statistiky jsou k~dispozici za poslední den, 3 dny, týden a~měsíc.
  \item[Contributors] Záložka poskytuje informace o~nejaktivnějších přispěvatelích, včetně automatizovaných účtů (botů). Zobrazuje počet commitů jak celkový, tak za každého přispěvatele zvlášť. Počet commitů je dynamicky agregován podle zvoleného časového období (minimálně 1 týden). Časový rozsah lze vybrat od posledních dvou týdnů po celou dobu existence repozitáře. Pokud repozitář obsahuje méně než 10~000 commitů, zobrazí tento pohled i~počet změněných řádků pro každého přispěvatele.
  \item[Commits] Záložka obsahuje graf celkového počtu commitů za poslední rok, přičemž počet commitů je agregován po týdnech.
\end{description}

Přehledný způsob zobrazení většího množství \PR nebo issues v~jednom grafu zde chybí. V~záložce Issues je však možné filtrovat pomocí pokročilých filtrů, které umožňují kombinovat podmínky pomocí booleovských operátorů a~závorek. Výsledkem je seznam issues nebo \PR s~celkovým počtem otevřených a~zavřených položek.

Dalším nástrojem jsou GitHub Projects, kde jsou issues k~jednotlivým funkcionalitám rozděleny do stavů definovaných vlastníkem repozitáře (například TODO nebo IN PROGRESS). Tento nástroj však neposkytuje žádné statistiky a~slouží spíše k~organizaci práce než k~získávání informací o~projektu.

\subsection{Zhodnocení}
\label{subsec:github_stats:evaluation}

Nativní statistiky GitHubu poskytují pouze základní přehled o~celkovém projektu. Jejich hlavní výhodou je nativní dostupnost bez nutnosti cokoliv nasazovat nebo udržovat.

Statistiky ohledně \PR nebo issues jsou zde obtížněji dostupné a~pro jejich získání je nutné sledovaný repozitář dobře znát. Na druhou stranu, k~nalezeným \PR nebo issues lze přejít přímo z~výsledků filtrování.

% **Jak GitHub kriticky zhodnotit (Zásadní pro tvou obhajobu):**

% * **Silné stránky (Plusy):** * Je to "out-of-the-box" řešení. Není potřeba nic instalovat, nasazovat ani udržovat.
%   * Data jsou 100% aktuální a UI je vývojářům důvěrně známé.
%   * Skvělé pro menší až střední projekty nebo pro letmý pohled.

% * **Slabé stránky / Omezení (Proč GitHub nestačí):**
%   * **Roztříštěnost informací:** Chybí jediný zastřešující "dashboard" na jednu obrazovku. Pokud chce manažer nebo vývojář zjistit zdraví projektu, musí proklikávat mezi záložkami Insights, Actions, Issues a Projects.
%   * **Omezená agregace napříč repozitáři:** Pokud se projekt skládá z 5 různých repozitářů, nativní UI GitHubu je neumí snadno spojit do jednoho uceleného grafu nebo metriky (což je u mikroservis častý problém).
%   * **Chybí kontextuální hloubka pro specifické workflow:** Nativní UI ti neřekne specifické metriky, jako např.: *"Jaká je průměrná doba čekání na code review u tohoto konkrétního týmu?"* nebo *"Který konkrétní test v CI pipelině selhává nejčastěji za poslední měsíc?"*.
